
\chapter[Design \textsuperscript{Kugler}]{Design}
\label{chap:Design}
The purpose of this chapter is to give an overview over the implementation design of Beat the Score. This chapter focuses on the communication between the different parts. 

\section{Data structures}

\subsection{{\tt Score}}

The data structure for songs. It consists of meta information about the song such as the title, the interprets and a list of tracks. This structure is used to store the songs from the GP5 files (see section \ref{sec:GP5}) with which the player input is compared to.

\subsection{{\tt Track}}

A {\tt Track} is a line of notes and information about the instrument that plays it. Additionally to the tracks in the {\tt Score}, we also store the player input in a {\tt Track}.

\subsection{{\tt Note}}

To store informations about individual notes. The most important ones are:
\begin{description}
	\item[time] in milliseconds after the start of the song. 
	\item[duration] in milliseconds.
	\item[pitch] as used in the MIDI specification.
\end{description}

\section{Program modules and communication}

For modularity the program has been split up into multiple parts:
\begin{description}
	\item[Inputs] \hfil \\ Handles inputs of an instrument.
	\item[Visuals] \hfil \\ Displays notes.
	\item[Sound] \hfil \\ Synthesizes an audio output.
	\item[Evaluator] \hfil \\ Compares the player input to the actual score and calculates how many points the player receives.
	\item[Main-Game] \hfil \\ Establishes and removes connections between the other parts when necessary. The GUI mostly communicates with this class.
\end{description}

Most communication between the different parts are done over Qt's Signal-/Slot-concept (see section \ref{sec:Signal-/Slot}).

\subsection{Inputs}
This part consists of the superclass {\tt Input} and an implementation for each input method. Scope of the project was only the support for MIDI keyboards over USB ({\tt MidiInput}) and guitars over phone connectors ({\tt AudioInput}, see chapter \ref{chap:Audio analysis}) but support for additional input methods can be added easily. 

\begin{sloppypar}
To create instances of those classes the class {\tt MidiInputManager} and {\tt AudioInputManager} are implemented. The purpose of those is to look for connected hardware of the respective type. 
\end{sloppypar}

{\tt Input} provides the following signals:
\begin{description}
	\item[{\tt noteOn(Note *note)}] \hfil \\ when a note is activated.
	\item[{\tt noteOff(Note *note)}] \hfil \\ when a note is released.
	\item[{\tt notePitchChanged(Note *note, float pitch)}] \hfil \\ when pitch bending is applied to a playing note.
\end{description}

\subsection{Visuals}
This part handles drawing the notes to the screen. It consists of the superclass {\tt Visuals} and its implementations {\tt PianoRoll}, {\tt ScoreView} and {\tt TablatureView}. The subclasses are detailed further in chapter \ref{chap:In-Game User Interface}. 

Additional display options can be added easily.

{\tt Visuals} provides the following slots:
\begin{description}
	\item[{\tt noteOn(Note *note)}] \hfil \\ when the player begins to play a note to make the note flash up in the view.
	\item[{\tt stateChanged(State state)}] \hfil \\ to react to pausing and resuming of the game in the trainings mode.
\end{description}

Further updates to player inputs are not required as an instance of {\tt Visuals} has a pointer to the {\tt Track} the {\tt Evaluator} writes into.

\subsection{Sound}
Responsible for audio output are the classes {\tt MidiNoteOutput} for the player input and {\tt NoteOutput} for the backing track. Both write into an instance of the class {\tt AudioBuffer} which is responsible for the mixing. By splitting it up like that we can avoid playing the notes from instruments that already have an audio output (in the scope of this project only guitars) by simply not connecting the {\tt Input} with the {\tt MidiNoteOutput}. How the synthesis of the audio output works is detailed in chapter \ref{chap:Audio output}.

The class {\tt MidiNoteOutput} has the following slots:
\begin{description}
	\item[{\tt noteOn(Note *note)}] \hfil \\ for when a note is activated.
	\item[{\tt noteOff(Note *note)}] \hfil \\ for when a note is released.
\end{description}

Those slots are compatible to the signals of the {\tt Input}s. Note that there is no corresponding slot for the {\tt notePitchChanged}-signal since we do not support pitch bending for MIDI instruments.

To play the backing tracks it stores a list of pointers to the corresponding {\tt Track} instances.

\subsection{Evaluator}
The linking piece between the inputs and the game. This part only consists of the class {\tt Evaluator} but could be extended by additional play modes in the future. 

The {\tt Evaluator} listens to notes with the slots 
\begin{description}
	\item[{\tt noteOn(Note *note)}] \hfil \\ for when a note is activated.
	\item[{\tt noteOff(Note *note)}] \hfil \\ for when a note is released.
	\item[{\tt notePitchChanged(Note *note, float pitch)}] \hfil \\ when pitch bending is applied to a playing note.
\end{description}
and reacts to them with the signal
\begin{description}
	\item[{\tt noteAdded(Note *note)}]\hfil \\ for further updates on the note, for example when the note is released, corresponding signals will be emitted by {\tt note} itself.
\end{description}

Transforming the stream of {\tt noteOn}- and {\tt noteOff}-signals into a stream of {\tt noteAdded}-signals with and observable {\tt Note} object has the advantage that we don't need to match the {\tt noteOff}-signals to the corresponding {\tt noteOn}-signals and also don't have to deal with notes being \emph{stuck} when the {\tt noteOff}-signal got lost because the object got disconnected from the input. Once the {\tt noteAdded} has been received by an object it can listen to the duration of the note which will be set as soon as the note is released. Even if the listening object got disconnected from the input, it can react to the end of the note properly.

The evaluator is also responsible for matching the player input with the actual score and award points depending on how close the note is to the actual score. The rules here are that the note must be in-key and the starting and ending points must be within a certain radius. If a new note is played which is closer to the actual note than a previously played note the closer one replaces the old on. 

\subsection{Main-Game}
This part consists of a single object of type {\tt MainGame} which is present throughout the whole runtime. It provides signals and properties accessible by the GUI to communicate the state of the game and the GUI calls its functions to control the game. This part only describes the user interactions from the \cpp side, for the QML side see chapter \ref{chap:UI}.

When the program is started the {\tt mainGame} object initializes the input managers and audio outputs and connects the {\tt MidiInput}s with the {\tt MidiNoteOutput} so the user can hear the notes they play already in the menu and to allow note-based navigation (see section \ref{sec:sound based navigation}).

The {\tt mainGame} object provides multiple functions to change the state of the game:
\begin{description}
	\item[{\tt songSelection()}] \hfil \\ to enter the song selection screen. It loads the default settings from the database so the GUI can preselect them.
	\item[{\tt startSong()}] \hfil \\ prepares a song for playing.
	\item[{\tt startAgain()}] \hfil \\ resets the user input so the song can be played again.
	\item[{\tt startReplay()}] \hfil \\ prepares playback of the previous user input.
\end{description}

Some of these states allow usage of additional functions. These are described in the following sections.

\subsubsection{{\tt songSelection()}}
The default settings are loaded from a local database (see section \ref{sec:Preferences}). If the user just played a song when this function is called all data of that song is released and the {\tt MidiInput}s are connected to the {\tt MidiNoteOutput} again (see section \ref{subsec:startSong}).

After this function has been called, multiple functions to change settings of the game can be used:
\begin{description}
	\item[{\tt selectSong(QString filename)}] \hfil \\ loads a song from the disk and prepares song-specific settings like track- and section-selection.
	\item[{\tt selectInput(int index)}] \hfil \\ selects the specified input device.
	\item[{\tt selectTrack(int index)}] \hfil \\ selects the specified track from the song.
	\item[{\tt setTempo(float factor)}] \hfil \\ changes the song to the specified tempo.
\end{description}

\subsubsection{{\tt startSong()}}
\label{subsec:startSong}
This function can be only called after all settings for a song have been set. It sets up the game in the following steps:
\begin{enumerate}
	\item All input devices that aren't used are disconnected. 
	\item A track for the user input is created.
	\item An evaluator is created and connected with the selected input device.
	\item The {\tt NoteOutput} gets a reference to the selected score.
\end{enumerate}

\subsubsection{{\tt startAgain()}}
After the completion of a song the player can choose to play again. To avoid destroying all objects just to create them again this is done with this extra function.

\subsubsection{{\tt startReplay()}}
Like the {\tt startAgain()} function this can be called after a song has been played. It will reset the visual but keep the player input. The inputs are not connected, so it's not possible to play new notes.